\section{The Appointed Texts}
\label{sec:appointed-texts}

The Latin below is that of the 1962 typical edition, \work{Missale Romanum}
(Typis Polyglottis Vaticanis), pp.~413--415, nos.~1699--1708, read on the page
images of the CMAA facsimile. Every element's wording already stands in the
Pustet \work{Missale Romanum} of Ratisbon, 1862 (pp.~354--355); the two books
differ only in spelling, accent, pointing, the form of the citations, the length
of the printed conclusions, and the additional orations that the 1862 book
prints under the rubrics of its day and the 1962 book no longer carries. The
English of the scriptural elements is the Douay--Rheims Bible in Bishop
Challoner's revision, which translates the Clementine Vulgate at the Missal's
own Vulgate numbering. It is given at the canonical verses, so that wherever a
chant departs from its verse the English and the Latin part company; each such
place is named beneath the text. Where the Douay--Rheims cannot carry the
Missal's form at all, as in the Offertory and the Communion, the English is that
of the anonymous translation printed in \work{The Roman Missal, Translated into
the English Language for the Use of the Laity} (Philadelphia: Eugene Cummiskey,
1861), which also supplies the English of the three orations. Neither English
is an approved liturgical translation of the 1962 Missal. Psalms are cited in
the Missal's numbering; the Hebrew numbering follows in parentheses.

\subsection{Introit}
\label{proper-introit}

\properlocus{\latin{Antiphona ad Introitum}, printed citation \latin{Esth.\ 13, 9 et 10-11}; verse Ps 118:1 (Heb.\ 119). Missal p.~413, no.~1699.}

\begin{latinproper}
In voluntáte tua, Dómine, univérsa sunt pósita, et non est qui possit
resístere voluntáti tuæ: tu enim fecísti ómnia, cælum et terram, et univérsa
quæ cæli ámbitu continéntur: Dóminus universórum tu es. \textnormal{Ps.\ 118, 1}
Beáti immaculáti in via: qui ámbulant in lege Dómini. ℣.~Glória Patri.
\end{latinproper}

\begin{englishwitness}{Douay--Rheims (Challoner), Esth 13:9--11, the source of the antiphon}
And said: O Lord, Lord, almighty king, for all things are in thy power, and
there is none that can resist thy will, if thou determine to save Israel. Thou
hast made heaven and earth and all things that are under the cope of heaven.
Thou art Lord of all, and there is none that can resist thy majesty.
\end{englishwitness}

\begin{englishwitness}{Douay--Rheims (Challoner), Ps 118:1, the verse}
Blessed are the undefiled in the way, who walk in the law of the Lord.
\end{englishwitness}

\noindent The antiphon is not a continuous text of the Vulgate, and the
Douay--Rheims above translates Mardochai's prayer as the Clementine Vulgate gives
it, not the Missal's antiphon. The Vulgate reads \latin{in ditione enim tua
cuncta sunt posita}; the Missal begins \latin{In voluntate tua \dots\ universa
sunt posita}, so that the will which no one resists is named twice, in the
first clause and in the second. The Missal puts \latin{resistere voluntati tuae}
in another order, adds \latin{enim} and \latin{omnia} (\latin{tu enim fecisti
omnia}), reads \latin{universa quae \dots\ continentur} for \latin{quidquid
\dots\ continetur} and \latin{Dominus universorum tu es} for \latin{Dominus
omnium es}. It leaves out the vocative \latin{Domine, Domine rex omnipotens}, the
condition \latin{si decreveris salvare Israel}, ``if thou determine to save
Israel'', and the closing \latin{nec est qui resistat maiestati tuae}. The 1861
hand missal renders the antiphon's opening sentence ``All things, O Lord, are
in thy power: and no one can resist thy will.'' The Vulgate counts the psalm's
title, \latin{Alleluia}, and the letter \latin{Aleph} inside its first verse;
the Missal's verse begins after them. After the \latin{Gloria Patri} the
antiphon is repeated (RGMR 427).

\subsection{Collect}
\label{proper-collect}

\properlocus{\latin{Oratio}. Missal p.~414, no.~1700.}

\begin{latinproper}
Famíliam tuam, quǽsumus, Dómine, contínua pietáte custódi: ut a cunctis
adversitátibus, te protegénte, sit líbera; et in bonis áctibus tuo nómini sit
devóta. Per Dóminum nostrum Iesum Christum, Fílium tuum: Qui tecum vivit et
regnat in unitáte Spíritus Sancti Deus: per ómnia sǽcula sæculórum. R.~Amen.
\end{latinproper}

\begin{englishwitness}{Cummiskey, The Roman Missal (1861), XXI Sunday after Pentecost, Collect}
Preserve thy family, O Lord, we beseech thee, by thy constant mercy: that under
thy protection, it may be freed from all adversity; and by the practice of good
works, continue devoted to thy name. Thro'.
\end{englishwitness}

\noindent ``By thy constant mercy'' stands for \latin{continua pietate}, the
unbroken \latin{pietas} of God toward his household; ``by the practice of good
works'' for \latin{in bonis actibus}. The Missal prints this conclusion in full,
with the response sign before \latin{Amen}, here written R. On 18~October 2026
the oration \latin{pro Fidei propagatione} is joined to the Collect under its
one conclusion in every Mass, and in a low or conventual Mass the Collect of St
Luke follows under its own (below).

\subsection{Epistle}
\label{proper-epistle}

\properlocus{\latin{Léctio Epístolæ beáti Pauli Apóstoli ad Ephésios}. Eph 6:10--17. Missal p.~414, no.~1701.}

\begin{latinproper}
Fratres: Confortámini in Dómino, et in poténtia virtútis eius. Indúite vos
armatúram Dei, ut possítis stare advérsus insídias diáboli. Quóniam non est
nobis colluctátio advérsus carnem et sánguinem: sed advérsus príncipes, et
potestátes, advérsus mundi rectóres tenebrárum harum, contra spirituália
nequítiæ, in cæléstibus. Proptérea accípite armatúram Dei, ut possítis
resístere in die malo, et in ómnibus perfécti stare. State ergo succíncti lumbos
vestros in veritáte, et indúti lorícam iustítiæ, et calceáti pedes in
præparatióne Evangélii pacis: in ómnibus suméntes scutum fidei, in quo possítis
ómnia tela nequíssimi ígnea exstínguere: et gáleam salútis assúmite: et gládium
spíritus, quod est verbum Dei.
\end{latinproper}

\begin{englishwitness}{Douay--Rheims (Challoner), Eph 6:10--17}
Finally, brethren, be strengthened in the Lord and in the might of his power.
Put you on the armour of God, that you may be able to stand against the deceits
of the devil. For our wrestling is not against flesh and blood; but against
principalities and powers, against the rulers of the world of this darkness,
against the spirits of wickedness in the high places. Therefore, take unto you
the armour of God, that you may be able to resist in the evil day and to stand
in all things perfect. Stand therefore, having your loins girt about with truth
and having on the breastplate of justice: And your feet shod with the
preparation of the gospel of peace. In all things taking the shield of faith,
wherewith you may be able to extinguish all the fiery darts of the most wicked
one. And take unto you the helmet of salvation and the sword of the Spirit
(which is the word of God).
\end{englishwitness}

\noindent The Missal opens the lesson with the bare address \latin{Fratres}
where the Vulgate's verse begins \latin{De cetero, fratres}, which the English
renders ``Finally, brethren''. The Missal spells \latin{exstinguere} for the Clementine's
\latin{extinguere}; every other word agrees with the Vulgate. The lesson stops
at verse 17. In the letter Paul goes on at once to ask for prayer ``at all times
in the spirit'' and for himself, ``for which I am an ambassador in a chain''
(6:18--20); the Mass does not read those verses.

\subsection{Gradual}
\label{proper-gradual}

\properlocus{\latin{Graduale}, Ps 89:1--2 (Heb.\ 90), respond and verse. Missal p.~414, no.~1702.}

\begin{latinproper}
Dómine, refúgium factus es nobis, a generatióne et progénie. ℣.~Priúsquam
montes fierent, aut formarétur terra et orbis: a sǽculo et usque in sǽculum tu
es, Deus.
\end{latinproper}

\begin{englishwitness}{Douay--Rheims (Challoner), Ps 89:1--2, without the title}
Lord, thou hast been our refuge from generation to generation. Before the
mountains were made, or the earth and the world was formed; from eternity and to
eternity thou art God.
\end{englishwitness}

\noindent The Missal reads \latin{a generatione et progenie}, ``from generation
and offspring'', where the Clementine reads \latin{a generatione in
generationem}; the English, which renders the Vulgate, cannot show the
difference, and the 1861 hand missal renders the Missal's respond with the same
words as the Douay. The Gradual leaves out the psalm's title, \latin{Oratio
Moysi, hominis Dei}, ``A prayer of Moses the man of God''. The verse agrees
with the Vulgate.

\subsection{Alleluia}
\label{proper-alleluia}

\properlocus{Ps 113:1 (Heb.\ 114:1). Missal p.~414, no.~1703.}

\begin{latinproper}
Allelúia, allelúia. ℣.~\textnormal{Ps.\ 113, 1} In éxitu Israël de Ægýpto,
domus Iacob de pópulo bárbaro. Allelúia.
\end{latinproper}

\begin{englishwitness}{Douay--Rheims (Challoner), Ps 113:1}
When Israel went out of Egypt, the house of Jacob from a barbarous people:
\end{englishwitness}

\noindent The verse agrees with the Vulgate, without its title
\latin{Alleluia}. The Vulgate's Ps~113 joins in one psalm the Hebrew Pss~114 and
115; the verse sung here is the first of the Exodus psalm that goes on to the
sea that fled and the Jordan turned back.

\subsection{Gospel}
\label{proper-gospel}

\properlocus{\latin{Sequéntia sancti Evangélii secúndum Matthǽum}. Mt 18:23--35. Missal p.~414, no.~1704.}

\begin{latinproper}
In illo témpore: Dixit Iesus discípulis suis parábolam hanc: Assimilátum est
regnum cælórum hómini regi, qui vóluit ratiónem pónere cum servis suis. Et cum
cœpísset ratiónem pónere, oblátus est ei unus, qui debébat ei decem mília
talénta. Cum autem non habéret unde rédderet, iussit eum dóminus eius venúmdari,
et uxórem eius, et filios, et ómnia quæ habébat, et reddi. Prócidens autem servus
ille, orábat eum, dicens: Patiéntiam habe in me, et ómnia reddam tibi. Misértus
autem dóminus servi illíus, dimísit eum, et débitum dimísit ei. Egréssus autem
servus ille, invénit unum de consérvis suis, qui debébat ei centum denários: et
tenens suffocábat eum, dicens: Redde quod debes. Et prócidens consérvus eius,
rogábat eum, dicens: Patiéntiam habe in me, et ómnia reddam tibi. Ille autem
nóluit: sed ábiit, et misit eum in cárcerem donec rédderet débitum. Vidéntes
autem consérvi eius quæ fiébant, contristáti sunt valde: et venérunt, et
narravérunt dómino suo ómnia quæ facta fúerant. Tunc vocávit illum dóminus suus:
et ait illi: Serve nequam, omne débitum dimísi tibi, quóniam rogásti me: nonne
ergo opórtuit et te miseréri consérvi tui, sicut et ego tui misértus sum? Et
irátus dóminus eius, trádidit eum tortóribus, quoadúsque rédderet univérsum
débitum. Sic et Pater meus cæléstis fáciet vobis, si non remiséritis
unusquísque fratri suo de córdibus vestris.
\end{latinproper}

\begin{englishwitness}{Douay--Rheims (Challoner), Mt 18:23--35}
Therefore is the kingdom of heaven likened to a king, who would take an account
of his servants. And when he had begun to take the account, one was brought to
him, that owed him ten thousand talents. And as he had not wherewith to pay it,
his lord commanded that he should be sold, and his wife and children, and all
that he had, and payment to be made. But that servant falling down, besought
him, saying: Have patience with me, and I will pay thee all. And the lord of
that servant being moved with pity, let him go and forgave him the debt. But
when that servant was gone out, he found one of his fellow-servants that owed
him an hundred pence: and laying hold of him, he throttled him, saying: Pay what
thou owest. And his fellow-servant falling down, besought him, saying: Have
patience with me, and I will pay thee all. And he would not: but went and cast
him into prison, till he paid the debt. Now his fellow servants seeing what was
done, were very much grieved, and they came, and told their lord all that was
done. Then his lord called him: and said to him: Thou wicked servant, I forgave
thee all the debt, because thou besoughtest me: Shouldst not thou then have had
compassion also on thy fellow servant, even as I had compassion on thee? And his
lord being angry, delivered him to the torturers until he paid all the debt. So
also shall my heavenly Father do to you, if you forgive not every one his
brother from your hearts.
\end{englishwitness}

\noindent The Missal opens with the liturgical formula \latin{In illo tempore},
names the speaker and his hearers, \latin{Dixit Iesus discipulis suis parabolam
hanc}, and so replaces the Vulgate's first word, \latin{Ideo}, ``Therefore''.
That word ties the parable to Peter's question, ``Lord, how often shall my
brother offend against me, and I forgive him? till seven times?'', and to the
answer, ``till seventy times seven times'' (18:21--22), which the Mass does not
read. The Missal reads \latin{suffocabat} where the Clementine reads
\latin{suffocavit} (v.~28), and spells \latin{milia} and \latin{venumdari};
``throttled'' renders either verb. The Gospel cross before the title is printed
in the Missal and not reproduced here. After the Gospel the Missal prints the cue
\latin{Credo}.

\subsection{Offertory}
\label{proper-offertory}

\properlocus{\latin{Ant.\ ad Offertorium}, printed citation \latin{Iob 1}. Missal p.~415, no.~1705.}

\begin{latinproper}
Vir erat in terra Hus, nómine Iob: simplex et rectus, ac timens Deum: quem
Satan pétiit, ut tentáret: et data est ei potéstas a Dómino in facultátes et in
carnem eius: perdidítque omnem substántiam ipsíus et filios: carnem quoque eius
gravi úlcere vulnerávit.
\end{latinproper}

\begin{englishwitness}{Cummiskey, The Roman Missal (1861), XXI Sunday after Pentecost, Offertory}
There was a man whose name was Job, and that man was simple and upright, and
fearing God: and Satan asked to tempt him; and power was given him by the Lord
over his possessions, and over his flesh: and he destroyed all his substance,
and his sons: and he wounded his flesh with a grievous ulcer.
\end{englishwitness}

\noindent The antiphon is a narrative in the chant's own words, not a verse of
Job. Its first clause shortens Job 1:1, which the Douay renders ``There was a man
in the land of Hus, whose name was Job, and that man was simple and upright, and
fearing God, and avoiding evil.'' What follows summarises the testing:
\latin{quem Satan petiit, ut tentaret} has no verbal counterpart in the book;
\latin{data est ei potestas \dots\ in facultates et in carnem eius} gathers the
two permissions, ``all that he hath is in thy hand'' (1:12) and ``he is in thy
hand, but yet save his life'' (2:6); \latin{perdiditque omnem substantiam ipsius
et filios} sums up the four messengers of 1:13--19; and \latin{carnem quoque eius
gravi ulcere vulneravit} stands for 2:7, ``struck Job with a very grievous
ulcer''. The 1861 English is that book's own rendering of the Missal's Latin;
it leaves out \latin{in terra Hus}. The Missal prints the antiphon alone.

\subsection{Secret}
\label{proper-secret}

\properlocus{\latin{Secreta}. Missal p.~415, no.~1706.}

\begin{latinproper}
Súscipe, Dómine, propítius hóstias: quibus et te placári voluísti, et nobis
salútem poténti pietáte restítui. Per Dóminum nostrum.
\end{latinproper}

\begin{englishwitness}{Cummiskey, The Roman Missal (1861), XXI Sunday after Pentecost, Secret}
Mercifully receive, O Lord, these offerings, by which thou art pleased to be
appeased, and restore us to salvation by thy powerful goodness. Thro'.
\end{englishwitness}

\noindent \latin{Quibus et te placari voluisti} says that God willed to be
appeased by these offerings, and \latin{et nobis salutem potenti pietate
restitui} that by them he willed salvation to be restored to us; the English
turns the second infinitive into a petition. ``Powerful goodness'' stands for
\latin{potenti pietate}, the same \latin{pietas} the Collect asked to be
continual. After the Secret the Missal prints the direction \latin{Præfatio de
Ss\~{m}a Trinitate}.

\subsection{Communion}
\label{proper-communion}

\properlocus{\latin{Antiphona ad Communionem}, Ps 118:81, 84 and 86 (Heb.\ 119). Missal p.~415, no.~1707.}

\begin{latinproper}
In salutári tuo ánima mea, et in verbum tuum sperávi: quando fácies de
persequéntibus me iudícium? iníqui persecúti sunt me, ádiuva me, Dómine Deus
meus.
\end{latinproper}

\begin{englishwitness}{Cummiskey, The Roman Missal (1861), XXI Sunday after Pentecost, Communion}
My soul hath looked to be saved by thee, and hath relied on thy word: when wilt
thou judge those who persecute me? The wicked ones have long persecuted me, help
me, O Lord my God.
\end{englishwitness}

\noindent The Communion sings three half-verses of the eleventh strophe of
Ps~118, the strophe \latin{Caph}, and changes each. Where the Vulgate reads
\latin{Defecit in salutare tuum anima mea: et in verbum tuum supersperavi}, the
Douay ``My soul hath fainted after thy salvation: and in thy word I have very
much hoped'', the Missal has no verb of fainting, \latin{In salutari tuo anima
mea}, and the simple \latin{speravi}. Of verse 84, ``How many are the days of thy
servant: when wilt thou execute judgment on them that persecute me?'', it sings
only the second question. Of verse 86, ``All thy statutes are truth: they have
persecuted me unjustly, do thou help me'', it sings only the second half, reads
\latin{iniqui}, ``the wicked'', for the Vulgate's adverb \latin{inique},
``unjustly'', and adds \latin{Domine Deus meus}. The 1861 English is that book's
own rendering of the Missal's Latin; its ``long'' has no word in the Latin.

\subsection{Postcommunion}
\label{proper-postcommunion}

\properlocus{\latin{Postcommunio}. Missal p.~415, no.~1708.}

\begin{latinproper}
Immortalitátis alimóniam consecúti, quǽsumus, Dómine: ut, quod ore percépimus,
pura mente sectémur. Per Dóminum nostrum.
\end{latinproper}

\begin{englishwitness}{Cummiskey, The Roman Missal (1861), XXI Sunday after Pentecost, Postcommunion}
Having received the food of immortality, we beseech thee, O Lord, that what we
have taken with our mouths we may receive with pure souls. Thro'.
\end{englishwitness}

\noindent The English renders \latin{sectemur} by ``receive''. The Latin verb
means to follow after or pursue: what the mouth has taken, the pure mind is to
go after. The prayer has the shape of a prayer the priest says after his
Communion in the Order of Mass, \latin{Quod ore sumpsimus} (\work{Ordo Missae},
no.~1130), which also sets what the mouth has taken beside a pure mind.

\subsection{The orations added on 18 October 2026}

On this date the rubrics add two sets of prayers to the Sunday's three orations;
neither belongs to the formulary. The oration \latin{pro Fidei propagatione}, its
Collect \latin{Deus, qui omnes homines vis salvos fieri}, Secret \latin{Protector
noster, aspice, Deus} and Postcommunion \latin{Redemptionis nostrae munere
vegetati} (printed in the Missal's \work{Tabula orationum}, p.~427,
nos.~1806--1808), is joined to each Sunday oration under one conclusion in every
Mass, unless the Ordinary has appointed another day for the missions (RGMR 453).
The commemoration of St Luke, his Collect \latin{Interveniat pro nobis},
Secret \latin{Donis caelestibus da nobis} and Postcommunion \latin{Praesta,
quaesumus, omnipotens Deus} (pp.~689--690, nos.~3883, 3889 and 3891), follows
under its own conclusion in a low or conventual Mass only (RG 108, 111;
RGMR 437~a).
